\documentclass[11pt,letterpaper]{article}
\usepackage{jahnel}

\runningtitle{Greedy paths in a Poisson line environment}
\hypersetup{pdftitle={Greedy paths in a Poisson line environment},pdfauthor={Christian Hirsch, Benedikt Jahnel}}

\title{\textbf{Greedy paths in a Poisson line environment\\ and the subcritical phase}\\[0.6em] \large Working note}

\author{%
  Christian Hirsch\thanks{Institut f\"ur Mathematische Stochastik, Georg-August-Universit\"at G\"ottingen, Goldschmidtstra\ss e~7, 37077 G\"ottingen, Germany. \texttt{christian.hirsch@uni-goettingen.de}. \orcid{0000-0003-4136-3740}.}
  \and
  Benedikt Jahnel\thanks{Institut f\"ur Mathematische Stochastik, Technische Universit\"at Braunschweig, Universit\"atsplatz~2, 38106 Braunschweig, Germany, and Weierstrass Institute for Applied Analysis and Stochastics, Mohrenstra\ss e~39, 10117 Berlin, Germany. \texttt{benedikt.jahnel@tu-braunschweig.de}. \orcid{0000-0002-4212-0065}.}%
}

\date{\today}

\begin{document}

\maketitle
\thispagestyle{plain}

\begin{abstract}
We examine the greedy-path approach to the subcritical phase of isotropic Poisson line Cox percolation. A quenched criterion gives a positive critical point intensity if the maximal sum of line counts along lattice paths grows linearly. We retain this condition as a conjecture and explain the dependence that prevents a direct application of the theory for independent weights. We correct the earlier moment calculation: finite paths have finite exponential moments, although those moments grow too rapidly for the proposed direct bound. We also withdraw the earlier polynomial lower bound for the dyadic run count. The corrected interval rate is of order $\mu/t^2$, and the accompanying calculation gives an exponential upper bound for the specified construction using a long pair of lines and individual additional intervals. Controlling connected runs and the total reward along paths remains open.
\end{abstract}

% =================================================================
\section{Setting and the greedy criterion}
\label{sec:setting}

Let $\cL$ be a Poisson process of lines in~$\R^2$ with intensity measure $(\mu/\pi)\dd\theta\dd s$ in the parametrisation $\ell(\theta,s)=\{x\colon x\cdot(\cos\theta,\sin\theta)=s\}$, $\theta\in[0,\pi)$. The lines meeting a convex body~$K$ have mean number $(\mu/\pi)$ times the perimeter of~$K$. The Boolean model $\Xi_\lambda$ is the union of disks of radius~$r$ around the points of a Poisson process with intensity~$\lambda$ per unit length on every line, and $\lambda_c(\mu,r)$ is its critical intensity~\cite{HJ26}. By scaling, $\lambda_c(\mu,r)=r^{-1}\lambda_c(\mu r,1)$, and we take $r=1$ throughout.

For $z\in\Z^2$ let $Q(z)\coloneqq z+[0,1]^2$, let $N(z)$ be the number of lines meeting~$Q(z)$, a Poisson variable with mean $4\mu/\pi$, and let $\operatorname{strip}(\ell)\coloneqq\{z\colon\ell\cap Q(z)\neq\varnothing\}$. A \emph{$2$-path} of length~$n$ is a self-avoiding sequence $\gamma=(z_0,\dots,z_n)$ in~$\Z^2$ with $\lVert z_i-z_{i-1}\rVert_\infty\le2$; $\Gamma_n(z_0)$ is the set of these paths, and $\#\Gamma_n(z_0)\le24^n$.

Define
\begin{equation*}
  Y_n(z_0)\coloneqq\sum_{\gamma\in\Gamma_n(z_0)}\ \prod_{i=0}^nN(z_i),\qquad g^*\coloneqq\limsup_{n\to\infty}\frac1n\log\max_{\lVert z_0\rVert_\infty\le n}Y_n(z_0)\,.
\end{equation*}

\begin{theorem}[{Quenched path criterion, \cite[Theorem~5.1]{HJ26b}}]
\label{thm:criterion}
For every locally finite~$\cL$, the quenched critical intensity satisfies $\lambda_c(\cL)\ge\e^{-g^*}/\sqrt2$. If $\P(g^*<\infty)>0$, then $\lambda_c(\mu,1)>0$.
\end{theorem}

The proof is a union bound over occupied $2$-paths: given the lines, boxes are occupied independently with probability at most $\sqrt2\,\lambda N(z)$, and an unbounded occupied component through~$Q(z_0)$ contains occupied $2$-paths of every length from~$z_0$. Positivity of $\lambda_c(\mu,1)$ then follows from the quenched zero--one law~\cite[Proposition~3.3]{HJ26b}. The criterion involves only the line process and imposes no restriction on~$\mu$.

\subsection{Linear reward}

The multiplicative path sum can be bounded using an additive reward. Counting the lines at each visited square gives a quantity that is linear in the line process. A finite growth rate for this reward is a sufficient condition for the quenched criterion. Define
\begin{equation}
\label{eq:h1}
  h_1^*\coloneqq\limsup_{n\to\infty}\frac1n\max_{\lVert z_0\rVert_\infty\le n}\ \max_{\gamma\in\Gamma_n(z_0)}R(\gamma),\qquad
  R(\gamma)\coloneqq\sum_{i=0}^nN(z_i)=\sum_{\ell\in\cL}\#\bigl(\gamma\cap\operatorname{strip}(\ell)\bigr)\,.
\end{equation}
Since $\log N\le N$ for $N\ge1$ and $\#\Gamma_n\le24^n$, we have $g^*\le\log24+h_1^*$. The quantity $h_1^*$ is translation invariant, hence almost surely constant by ergodicity of~$\cL$. The reward $R(\gamma)$ is linear in the line process: $R(\gamma)=\sum_\ell f_\gamma(\ell)$ with $f_\gamma(\ell)=\#(\gamma\cap\operatorname{strip}(\ell))$. Trivially $h_1^*\ge1$, since every box met by a given line has $N\ge1$.

The remaining sufficient estimate is the following assertion about the line process alone.

\begin{conjecture}[Linear greedy growth]
\label{conj:greedy}
For every $\mu>0$, almost surely $h_1^*<\infty$.
\end{conjecture}

By \cref{thm:criterion}, \cref{conj:greedy} implies $\lambda_c(\mu,1)\ge\e^{-\log24-h_1^*}/\sqrt2>0$. Scaling gives the conclusion for general $r>0$, with the line intensity replaced by $\mu r$ in the unit-radius model.

\subsection{Comparison with greedy lattice animals}

Independent weights provide a useful benchmark for this problem. Their theory controls the total reward collected by connected sets under a condition on the marginal tail. Here we use that theory to identify the extra issue created by shared lines.

Let $d\ge2$, and let $(W_z)_{z\in\Z^d}$ be independent and identically distributed nonnegative weights with common law that of $W$. For connected sets $\xi$ of $n$ sites containing the origin, define
\[
 M_n=\max_{\xi\ni0,\,|\xi|=n}\sum_{z\in\xi}W_z.
\]
Classical upper bounds were developed in~\cite{CGGK93}. The almost-sure linear growth of $M_n$ follows under the sufficient condition
\begin{equation}
\label{eq:martin}
  \int_0^\infty\P(W>w)^{1/d}\dd w<\infty
\end{equation}
by Martin~\cite{Mar02}. These results assume independence of the weights. The Poisson line field has unbounded dependence along every line, so its marginal moment bounds alone do not verify the hypotheses.

The field $N$ has superexponentially light one-point tails, $\P(N(z)\ge k)\le(4\e\mu/(\pi k))^k$, so the one-point condition~\eqref{eq:martin} holds with room to spare. The difficulty is the dependence: a line contributes to $N$ along its whole strip, and \emph{bundles} of nearly parallel lines create long rows of boxes with large~$N$. The following two facts quantify this dependence~\cite[Lemma~7.1]{HJ26b}.

\begin{lemma}[Shared lines]
\label{lem:shared}
For every $k\in\N$ there is $c_k<\infty$ such that, almost surely, for all large~$n$, any two boxes in $[-n,n]^2$ at distance at least $c_kn^{6/k}$ share fewer than~$k$ lines. Moreover, almost surely $\max_{\lVert z\rVert_\infty\le n}N(z)\le\log n$ for all large~$n$.
\end{lemma}

% =================================================================
\section{The direct moment and path-counting bounds}
\label{sec:firstmoment}

Rare collections of nearly parallel lines contribute to many paths at once. This causes rapid growth of the annealed moments used in the direct path bound. We first give the corrected finite-path estimate, then explain which counting obstruction survives the correction.

Let $\gamma_n=(0,e_1,\ldots,ne_1)$ and let $M$ be the number of lines meeting both vertical sides of $[0,n+1]\times[0,1]$. Its mean $\nu_n$ satisfies $\mu/(5\pi(n+1))\le\nu_n\le2\mu/\pi$, by~\cite[Lemma~2.1]{HJ26b}, and every such line meets every square of $\gamma_n$. Let $a_n=2\mu(n+2)/\pi$, the mean number of lines meeting the whole rectangle. The next proposition corrects the former claim that the exponential moment is infinite at a finite path length.

\begin{proposition}[Exponential moments of a straight path]
\label{prop:annealed}
For every $\mu>0$, $s>0$, and integer $n\ge0$,
\[
 \nu_n\big(\e^{s(n+1)}-1\big)
 \le\log\E\e^{sR(\gamma_n)}
 \le a_n\big(\e^{s(n+1)}-1\big)<\infty.
\]
Consequently, $\lim_{n\to\infty}n^{-1}\log\E\e^{sR(\gamma_n)}=\infty$.
\end{proposition}

\begin{proof}
The number $M$ is Poisson with mean $\nu_n$, and $R(\gamma_n)\ge(n+1)M$. If $M'$ is the number of lines meeting the enclosing rectangle, then $M'$ is Poisson with mean $a_n$ and $R(\gamma_n)\le(n+1)M'$. Applying the Poisson generating function to the two bounds proves the displayed inequalities. The lower bound on $\nu_n$ proves the limit. Every line contributes at most $n+1$ to the reward, so the exponential-formula integral cannot be summed over larger incidences as in the former proof.
\end{proof}

The superexponential annealed growth of $Y_n(0)$ also remains valid. With $j=\lceil n/\log n\rceil$ for $n\ge3$, the straight path gives
\[
 \E Y_n(0)\ge j^{n+1}\P(M\ge j)
 \ge j^{n+1}\e^{-\nu_n}\frac{\nu_n^j}{j!}.
\]
Taking logarithms yields $n\log n-n\log\log n-O(n)$ as a lower bound. Thus, $n^{-1}\log\E Y_n(0)\to\infty$, which prevents a useful bound of the required exponential order through this expectation.

For the path-counting statement, put $n=2p$ and fix an integer $H\ge1$. There are $2^p$ distinct paths on row zero whose $p$ increments are $e_1$ or $2e_1$, and every visited square lies in the rectangle above. On $\{M\ge H\}$ each of these paths has reward at least $H(p+1)$. Their expected number with this reward is therefore at least $2^p\e^{-\nu_{2p}}\nu_{2p}^H/H!$. This construction corrects the former use of row $-1$, which is not traversed by the chosen lines. It shows why counting all these paths separately is inefficient; it does not exclude every first-moment argument after a different decomposition.

\section{The sandwich identity}
\label{sec:sandwich}

The level sets $D_m\coloneqq\{z\colon N(z)\ge m\}$ give $R(\gamma)=\sum_{m\ge1}\#(\gamma\cap D_m)$, and $\#(\gamma\cap D_m)\le\#\gamma$ for $m\le2$. A symmetric enumeration of $D_m$ by $m$-subsets of lines overcounts a $j$-bundle by a factor $\binom jm$, which destroys the exponential tails needed later (Bell-number blow-up). The following exact identity avoids this.

For a box~$z$ with $N(z)\ge3$, order the lines through~$Q(z)$ by their angle in~$[0,\pi)$. A line~$\ell$ is \emph{sandwiched at~$z$} if it meets~$Q(z)$ and is neither the smallest nor the largest angle among these lines. Let $Z(\ell)\coloneqq\{z\colon\ell\text{ sandwiched at }z\}\subseteq\operatorname{strip}(\ell)$.

\begin{lemma}[Sandwich identity]
\label{lem:sandwich}
For every box~$z$, $(N(z)-2)^+=\#\{\ell\colon\ell\text{ sandwiched at }z\}$. Consequently, for every path~$\gamma$,
\begin{equation}
\label{eq:sandwich}
  R(\gamma)\le2\,\#\gamma+\sum_{\ell\in\cL}\#\bigl(\gamma\cap Z(\ell)\bigr)\,.
\end{equation}
Replacing "neither the smallest nor the largest" by "neither among the $a$ smallest nor among the $b$ largest" gives $(N-a-b)^+$.
\end{lemma}

\begin{proof}
Among $N\ge3$ lines ordered by angle, exactly $N-2$ are neither first nor last. Summing over $z\in\gamma$ and interchanging the sums gives~\eqref{eq:sandwich}.
\end{proof}

On a set of $q$ squares each met by exactly $j\ge3$ lines, the excess counted by the identity is $(j-2)q$. The identity reduces \cref{conj:greedy} to bounding the time a path can spend in sandwiched boxes of the lines.

For a deterministic line $\ell$, each other nonparallel line produces the interval $J(\ell,\ell')$ consisting of points of $\ell$ within distance $2$ of $\ell'$. At relative direction angle $\alpha\in(0,\pi)$ its length is $4/\sin\alpha$, and crossings have intensity $(\mu/\pi)\sin\alpha\,\dd x\dd\alpha$ in arclength and angle coordinates. Consequently, intervals of length in $[t,2t]$ start at rate of order $\mu/t^2$, while the expected number of these intervals covering a fixed point is of order $\mu/t$. The accompanying note \texttt{corrected-geometry.pdf}, Lemma~2.1, gives exact formulas and includes intervals entering from outside a segment.

The lattice set $Z(\ell)$ is not literally a subset of the line. If two lines meet a unit square, each point of either line inside the square is within distance $\sqrt2<2$ of the other. Hence, the union of $Q(z)\cap\ell$ over $z\in Z(\ell)$ is contained in the intersection of two unions of intervals $J(\ell,\ell')$, using lines of smaller and larger normal angles, respectively. The reverse inclusion is unavailable. Connected components of these unions may concatenate many shorter intervals, so their length distribution requires a separate argument.

\section{Geometric density heuristics}
\label{sec:bundles}

The earlier density calculation provides a heuristic for a scale decomposition. It concerns finite collections of almost parallel lines close to a prescribed point. Passing from that calculation to a stationary collection of distinct connected runs requires definitions and bounds that have not been established.

For fixed integer $j\ge2$, lines meeting a fixed unit disk have measure of order $\mu$. Given one line, the measure of another meeting this disk at an angle of order $1/t$ from it is of order $\mu/t$. Thus, the expected count of the specified $j$-tuples has order $\mu^jt^{1-j}$, with constants depending on $j$. Dividing by the length $t$ suggests a density of order $\mu^jt^{-j}$ for distinct geometric stretches. This division is a heuristic until a rule for choosing and measuring such stretches has been given.

In particular, this computation does not prove that maximal runs have density $\mu^3t^{-3}$. It also does not show that distant stretches use disjoint lines: \cref{lem:shared} bounds the number of shared lines and permits a positive number of them. The former refinement $h_1^*=2+O(\mu^{3/2})$ and the series used to support it are withdrawn from the conjecture, since neither the necessary component bounds nor constants uniform in the number of lines have been obtained.

\section{Reassessment of the dyadic estimate}
\label{sec:dyadic}

The earlier scheme used the tentative density $\mu^3t^{-3}$ to select boxes of side $w=(t/\mu)^{3/2}$. It then claimed that a pair of close lines spanning a row of $m$ boxes creates more than $Km$ additional intervals of length comparable to $t$ with polynomial probability. The corrected calculation gives an exponential upper bound for this particular mechanism at sufficiently large $t$.

The measure of the first lines crossing a long row is of order $\mu w/m$. For a fixed first line, the measure of a second line staying within a fixed distance of it over length $mw$ is of order $\mu/(mw)$. Their product is therefore of order $\mu^2/m^2$, without the extra factor $w$ in the former proof. Given a deterministic pair, the remaining lines have their original Poisson law, and the expected number of individual intervals in the required length range along the row is of order
\[
 \frac{\mu mw}{t^2}=\frac{m}{\sqrt{\mu t}}.
\]
The endpoint correction for intervals entering from outside the row is of order $\mu/t$. Both contributions are computed exactly in \texttt{corrected-geometry.pdf}.

Proposition~1.1 of that accompanying note makes the geometry precise and applies the two-point Mecke formula to sum over candidate pairs. For each fixed $\mu>0$ and integer $K\ge1$, it proves an upper bound $C\mu^2m^{-2}2^{-Km}$ once $t$ is sufficiently large. Thus, the polynomial lower bound formerly called Proposition~5.1 is withdrawn. This correction does not prove an exponential bound for the full count of maximal runs, since those runs may arise from chains of shorter intervals or from intersections of longer intervals.

There are two further gaps in the proposed synthesis. The sum of the terms $2^k$ coming from the additive one in the number of boxes touched diverges unless large scales are treated separately; the previous $o(n)$ remainder was unsupported. Moreover, bounds for connected sets of boxes must specify the allowed sizes and locations before a union bound is applied. These issues remain even after the count for one pair has been corrected.

\section{Remaining problem}
\label{sec:outlook}

The unresolved goal remains \cref{conj:greedy}. The exact counting identity retains the excess over two lines per visited square, and the corrected pair estimate controls one source of that excess. A full proof must bound the reward collected by paths that pass between many different configurations of lines.

A useful next input would control connected components of the intersections of the two interval unions on a line, with multiplicities sufficient to recover the lattice reward. The estimate would then have to remain valid when these lines are chosen and visited by a path. Shared lines prevent replacing the resulting collection by an independent Poisson process without further argument, and the Mecke formula only removes the selection bias for the finite collection of lines explicitly integrated out.

The rigorous conclusions retained here are the quoted quenched criterion, its reduction to linear reward, the shared-line bounds, the corrected moments and path counts, and the counting identity. The additional pair estimate is proved in full in the accompanying correction note. Positivity of the critical intensity, the requisite component bounds, and a complete dyadic synthesis remain unproved. No conclusion that all possible approaches fail follows from these calculations.

\section*{Acknowledgements}

BJ acknowledges the financial support of the Leibniz Association within the Leibniz Junior Research Group on Probabilistic Methods for Dynamic Communication Networks (DYCOMNET) as part of the Leibniz Competition (grant no.~J105/2020), of the Berlin Mathematics Research Center MATH+ through the project EF-MA-Sys-2 on Information Flow \& Emergent Behavior in Complex Networks, and of the Deutsche Forschungsgemeinschaft (DFG) through Project P27 within the SPP~2265.

\section*{Use of AI tools}

The simulations, proofs, computations, figures and the text of this note were developed with substantial assistance of the AI systems Claude (Anthropic) and ChatGPT Astra~6 (OpenAI), in an interactive process directed by the authors.

\bibliographystyle{amsalpha}
\bibliography{references}

\end{document}
